\documentclass[11pt]{article}
\usepackage[margin=1in]{geometry}
\usepackage{amsmath,amssymb,booktabs,longtable,array}
\usepackage[hidelinks]{hyperref}
\newcommand{\dd}{\mathrm d}
\newcommand{\DE}{\mathrm{DE}}
\newcommand{\Pl}{\mathrm{Pl}}
\newcommand{\Mb}{M_b}
\title{Cold energy, dark energy, and a particle-physics bridge\\
An equation inventory and three executed closure discriminators}
\author{Zimmerman programme: research checkpoint CPB-2026-10-16}
\date{16 October 2026}
\begin{document}
\maketitle
\begin{abstract}
The physical theory remains open. This report separates vacuum stress,
mass-carrying excitations, and collective response, and records the equations
needed to connect them. Three routes are executed: a charge-preserving Higgs
production extension, an exact cold--vacuum exchange consistency relation,
and an exterior-mass lensing test. The simplest portal extension encounters
mass-potential and kinetic-evolution obstructions at the adopted benchmark;
a source-quartic extension repairs the global classical potential exactly;
the exchange relation supplies a falsifiable joint constraint; positive
spherical exterior mass cannot increase interior excess tangential shear.
These are conditional reductions and applications of known mechanisms,
not a particle discovery, a Standard Model derivation, or global novelty claims.
\end{abstract}

\section{Target, evidence, and conventions}
The requested goal is one admitted physical model producing the galactic law,
its finite mass supply, formation and source selection, the lensing metric,
and cosmological abundance and perturbations. A predictive effective theory
may contain measured constants. The stronger goal of deriving every constant
also requires a mechanism selecting those constants; it is a separate standard.
Neither standard is met by combining successful tests of different actions.

The baseline is Git revision
\texttt{370dfb58006e330f210f5e306297e822188e37ba} in a shared dirty checkout.
Inputs to the calculation are hashed in \texttt{run\_003/manifest.json}.
The controlling recent work order is
\path{sol61_push/THEORY_NEXT_STEPS_2026-10-15.md}, including its October 15
fourth-round successor. The frozen October 14 candidate is preserved.
Older README and retraction prose describe different stages and do not
override the dated candidate-specific results.

\textbf{MODEL} denotes a chosen action/state; \textbf{DERIVED} means conditional
on the displayed hypotheses; \textbf{TARGET} denotes a relation to explain;
\textbf{OPEN} denotes a missing implication. Equations from distinct models
below are not silently parts of one action. Microscopic formulas use
$\hbar=c=k_B=1$, signature $(+---)$; cosmological SI equations explicitly
use mass-equivalent densities $\rho_i$ and $p_i/c^2$. The weak lensing metric
uses $(-+++)$ separately. These conventions must be translated in a future
joined action. $\Phi_c$ below is the complex carrier; $\phi_g$ is gravitational
potential; $\mathbf P$ is polarization; $p$ is pressure.

The two acceleration footings remain separate:
\begin{equation}
 a_{0,\mathrm{can}}=9.3603\,10^{-11}\ \mathrm{m\,s^{-2}},\qquad
 a_{0,\mathrm{alt}}=1.1312\,10^{-10}\ \mathrm{m\,s^{-2}}.
\end{equation}
These are the frozen candidate's values; nearby older documents use slightly
different rounded values. The alternative total-density footing must not
be identified with vacuum density in the equations below. The new microscopic
controls do not choose or combine these galaxy footings.

\section{What cold and dark energy mean}
Cold energy is energy mostly stored as rest mass, with
$|p_c|\ll\rho_c c^2$ in the relevant frame. It is not an identified species.
Dark energy is a component with sufficiently negative pressure to accelerate
expansion. A constant positive vacuum has
\begin{equation}
 \epsilon_v=\rho_v c^2>0,\qquad p_v=-\epsilon_v,\qquad w_v=-1.
\end{equation}
A nonrelativistic isotropic distribution instead has, at leading order,
\begin{equation}
 \rho_c c^2\simeq mn c^2+\tfrac12 mn\langle v^2\rangle,
 \qquad p_c\simeq\tfrac13 mn\langle v^2\rangle.
\end{equation}
Low mean pressure does not establish a small perturbation sound speed,
short free-streaming length, correct abundance, or a galactic attractor.

A simple example of one field carrying both roles is the stipulated potential
\begin{equation}
 V(\Phi_c)=V_0+m^2|\Phi_c|^2+\lambda|\Phi_c|^4,\quad
 \Phi_c\simeq e^{-imt}\psi/\sqrt{2m},\quad n=|\psi|^2.
\end{equation}
In its dilute rotating, weak-interaction limit,
\begin{equation}
 \epsilon\simeq V_0+mn+\frac{\lambda n^2}{4m^2},\qquad
 p\simeq-V_0+\frac{\lambda n^2}{4m^2}.
\end{equation}
Here $mn$ is rest-energy density in natural units. The interaction is
repulsive for $\lambda>0$. The approximate adiabatic sound speed of the
nonvacuum part is
\begin{equation}
 c_{s,c}^2\simeq
 \frac{\lambda n/(2m^2)}{m+\lambda n/(2m^2)}.
\end{equation}
This illustrates coexistence, not the desired galaxy force. Changing $V_0$
leaves the flat-space scalar equation unchanged while changing vacuum gravity.
Thus writing one field does not fix the cold abundance, vacuum magnitude,
or $a_0$. Previously tested linked-square potentials also have excessive
pressure/scale-separation problems; their failures are not erased here.

\section{Dark energy, expansion, and the acceleration scale}
For FLRW geometry,
\begin{align}
 H^2&=\frac{8\pi G}{3}\sum_i\rho_i-\frac{k c^2}{a^2},\\
 \frac{\ddot a}{a}&=-\frac{4\pi G}{3}\sum_i(\rho_i+3p_i/c^2).
\end{align}
Vacuum may be written as density on the right or as $\Lambda$ on the left
of Einstein's equation, once:
\begin{equation}
 \rho_v=\frac{\Lambda c^2}{8\pi G},\quad
 H_\Lambda^2=\frac{\Lambda c^2}{3},\quad
 a_0=\kappa c\sqrt{G\rho_v}
 =\kappa c^2\sqrt{\frac{\Lambda}{8\pi}}.
\label{eq:scale}
\end{equation}
The last link is the programme's \textbf{TARGET}, with $\kappa=1/2$ fitted,
not derived. Consequently
\begin{equation}
 a_0^2=\frac{\Lambda c^4}{32\pi},\quad
 \epsilon_v=\frac{4a_0^2}{G},\quad
 \frac{a_0}{cH_\Lambda}=\frac{1}{2}\sqrt{\frac{3}{8\pi}}.
\end{equation}
These three are algebraic restatements of the adopted coefficient.
They are not three independent derivations. A constant offset of a response
potential remains a free vacuum term unless an actual mechanism fixes it.

For separately conserved dark energy,
\begin{equation}
 \frac{\rho_\DE(z)}{\rho_{\DE0}}
 =\exp\!\left[3\int_0^z\frac{1+w(z')}{1+z'}\dd z'\right],\qquad
 \frac{a_0(z)}{a_0(0)}=\sqrt{\frac{\rho_\DE(z)}{\rho_{\DE0}}}.
\end{equation}
For CPL, $w(a)=w_0+w_a(1-a)$,
\begin{equation}
 F(z)=(1+z)^{3(1+w_0+w_a)}e^{-3w_a z/(1+z)}.
\end{equation}
Constant vacuum gives constant $a_0$. Evolving CPL density, an $H(z)$ law,
and an event-horizon law are distinct hypotheses. An event-horizon model
instead requires a specified future expansion history:
\begin{equation}
 R_E(a)=ca\int_a^\infty\frac{\dd a'}{a'^2 H(a')},\qquad
 a_0(a)/a_0(1)=R_E(1)/R_E(a),
\end{equation}
when the integral exists. No existing seal is changed here.

\section{The galaxy law and the static polarization construction}
For spherical baryons,
\begin{equation}
 g_b=\frac{G\Mb(<r)}{r^2},\quad y=g_b/a_0,\quad
 g=\nu(y)g_b,\quad g\sim\sqrt{a_0g_b},\quad V_f^4=G\Mb a_0.
\end{equation}
Three distinct interpolations appearing in the corpus are
\begin{align}
 \nu_s(y)&=\tfrac12+\sqrt{\tfrac14+1/y},
 &g_b&=\frac{g^2}{g+a_0},\\
 \nu_q(y)&=\sqrt{1+1/y},
 &g^2-g_b^2&=a_0g_b,\\
 \nu_e(y)&=(1-e^{-\sqrt y})^{-1}.
\end{align}
They share a deep limit and differ elsewhere. They must not be called
equivalent. The October 14 static candidate uses $\nu_s$.

Put $K=4\pi G$, $u=K|\mathbf P|<a_0$, and
$\gamma=K L_*^2/2$. Its chosen static action density is
\begin{equation}
 \mathcal L_{S2}=-\frac{|\nabla\phi_g-K\mathbf P|^2}{2K}
 -\rho_b\phi_g-W_{\rm rem}(u)-\frac\gamma2|\nabla\mathbf P|^2,
\end{equation}
where the critical quadratic part has already been included in the square:
\begin{equation}
 W_{\rm rem}(u)=\frac1K\left[-\frac{u^2}{2}-a_0u
 -a_0^2\ln(1-u/a_0)\right],\qquad
 K W'_{\rm rem}=\frac{u^2}{a_0-u}.
\end{equation}
Variation, before a phase gate is added, gives
\begin{align}
 \nabla\cdot(\nabla\phi_g-K\mathbf P)&=K\rho_b,\\
 \nabla\phi_g-K\mathbf P-KW'_{\rm rem}(u)\widehat{\mathbf P}
 +\gamma\nabla^2\mathbf P&=0.
\end{align}
Outside stiffness-dominated regions, in the spherical aligned branch,
\begin{equation}
 g-u=g_b=\frac{u^2}{a_0-u},\qquad
 u=\frac{a_0g}{a_0+g},\qquad g_b=\frac{g^2}{g+a_0}.
\end{equation}
This is \textbf{DERIVED} from a potential chosen to give the law. It does not
derive that potential from microscopic physics. $L_*$ is a bounded input.
The static action alone supplies no inertia, collision operator, physical
clock, or condensed-baryon ownership prescription. A local AQUAL-like
reduction also cannot by itself establish absence of an external-field effect.

\section{Mass supply, currents, and free matter}
If $\mathbf P$ is actual cold-mass displacement, the proposed identification is
\begin{equation}
 \nabla\cdot\mathbf P=\rho_c-\bar\rho_c,\qquad
 \dot{\mathbf P}=-\mathbf j_c.
\end{equation}
Continuity alone implies only
$\nabla\cdot(\mathbf j_c+\dot{\mathbf P})=0$.
For homogeneous positive $\rho_0$, ordinary material kinetic energy requires
\begin{equation}
 T\ge\frac{|\mathbf j_c|^2}{2\rho_0},\qquad
 T_k\ge\frac{|\dot{\delta\rho}_k|^2}{2\rho_0k^2},\qquad
 I_{\rm mass}=\rho_0^{-1},\quad
 c_{\rm bare}^2=\gamma\rho_0=2\pi G\rho_0L_*^2.
\end{equation}
This reproduces the October 15 mass-response audit. Choosing an independent
fast response field changes the model; it does not remove mass transport.
For a localized contrast,
\begin{equation}
 \int\delta\rho\,\dd^3x=\oint\mathbf P\cdot\dd\mathbf S.
\end{equation}
Zero boundary flux implies compensation, not creation of positive mass
from an empty isolated domain.

In the zero-stiffness point-baryon $\nu_s$ branch put
$r_M=\sqrt{G\Mb/a_0}$ and $x=r/r_M$. Then
\begin{equation}
 \frac{M_t(<r)}{\Mb}=\frac{\sqrt{1+4x^2}-1}{2},\quad
 M_t=\frac{r^2(g-g_b)}{G}.
\end{equation}
An allocated stock $S_c=\mathcal R\Mb$ therefore gives
\begin{equation}
 x\le\sqrt{\mathcal R(\mathcal R+1)}=5.8426446\ldots
 \quad(\mathcal R=5.364).
\end{equation}
The numerical stock is an input inferred in a cosmological model, not a
derivation or a measured species identity. More generally
\begin{equation}
 M_{c,\rm target}=(\nu-1)\Mb,\quad
 \rho_{c,\rm target}=\frac{1}{4\pi r^2}\frac{\dd M_{c,\rm target}}{\dd r}.
\end{equation}
Both enclosed positivity and density positivity must hold.
The earlier bookkeeping $M_{\rm tot}=\nu\Mb+k M_{\rm hot}$ has fitted
$k\simeq3$; it is not a physical partition law. The newer work order retains
the failure of a gas-tracing free pool and requires a counted free population.

For a conditional extension where free mass supplies ordinary source
$A=g_b+g_f$ to the same responding phase,
\begin{equation}
 A=\frac{g^2}{g+a_0},\quad a_{\rm app}=\frac{g(g-g_b)}{g_b},\quad
 M_f(<r)=\Mb(<r)\frac{a_{\rm app}-a_0}{g+a_0}.
\end{equation}
The October 15 inverse diagnostic requires $M_f\ge0$ and $M'_f\ge0$.
Unchanged exact baryon-only adherence then leaves no such free mass inside
the radius. A different source/phase assignment requires different equations.

\section{The microscopic cold carrier}
The independent charge-protected scalar model is
\begin{align}
 \mathcal L_c&=|\partial\Phi_c|^2+|\partial X|^2-V_c,\\
 V_c&=m^2|\Phi_c|^2+4m^2|X|^2+
 \left|\sqrt{\Delta_c}X+\frac{g_c}{\sqrt{\Delta_c}}\Phi_c^2\right|^2
 +\lambda_{\rm rep}|\Phi_c|^4,\\
 M&=2m\sqrt{1+u_*^2},\quad \Delta_c=4m^2u_*^2,
 \quad\lambda_{\rm tot}=g_c^2/\Delta_c+\lambda_{\rm rep}.
\end{align}
Here $m,\Delta_c,\lambda_{\rm rep}>0$ and $g_c$ is real. Its equations are
\begin{equation}
 (\Box+m^2)\Phi_c+2g_c X\Phi_c^*
 +2\lambda_{\rm tot}|\Phi_c|^2\Phi_c=0,\qquad
 (\Box+M^2)X+g_c\Phi_c^2=0.
\end{equation}
The symmetry $(\Phi_c,X)\mapsto(e^{i\theta}\Phi_c,e^{2i\theta}X)$ yields
\begin{equation}
 j_Q^\mu=i(\Phi_c^*\partial^\mu\Phi_c-\Phi_c\partial^\mu\Phi_c^*)
 +2i(X^*\partial^\mu X-X\partial^\mu X^*),\quad \partial_\mu j_Q^\mu=0.
\end{equation}
In the isolated classical model $E\ge m|Q|$. This carrier-only bound does
not automatically include gravity, a portal, or a finite-temperature medium.

At a fixed regulator in the nonrelativistic approximation put
$\delta=M-2m$, $c_3=g_c/(2m\sqrt{2M})$, $k_c=\lambda_{\rm tot}/(4m^2)$.
With source annihilator $b$,
\begin{align}
 H&=mQ+D_c,\quad Q=\int(\psi^\dagger\psi+2b^\dagger b)\dd^3x,\\
 D_c&=T+\delta\int\left|b+\frac{c_3}{\delta}\psi^2\right|^2\dd^3x
 +\left(k_c-\frac{c_3^2}{\delta}\right)\int\psi^{\dagger2}\psi^2\dd^3x,\\
 k_c-\frac{c_3^2}{\delta}
 &=\frac{\lambda_{\rm rep}}{4m^2}
 +\frac{g_c^2}{8m^2M(M+2m)}>0.
\end{align}
Thus $0\le\langle T\rangle\le\langle D_c\rangle$. For initially resting
sources $D_c/(mQ)=\delta/(2m)\simeq u_*^2/2$.
The regulated homogeneous pressure bound uses
\begin{equation}
 r_c^2=\frac{c_3^2}{\delta k_c}<1,\quad
 b_\pm=\tfrac12(1\pm(1-r_c^2)^{-1/2}),\quad
 \frac{|p_c|}{\epsilon_c}\le\frac{b_+d}{1+d},\quad
 d=\frac{\langle D_c\rangle}{m\langle Q\rangle}.
\end{equation}
It is not a continuum quantum-field or perturbation theorem.

A resting source decays to two same-charge daughters, with
\begin{equation}
 p_*=mu_*,\quad v_*^2=\frac{u_*^2}{1+u_*^2},\quad
 \frac{p_{\rm free}}{\epsilon_{\rm free}}=\frac{u_*^2}{3(1+u_*^2)},\quad
 \Gamma_c=\frac{g_c^2}{8\pi M}\sqrt{1-\frac{4m^2}{M^2}}.
\end{equation}
At threshold and away from an unresolved resonance,
\begin{align}
 \sigma_0&=\frac{\lambda_{\rm rep}^2}{8\pi m^2},\quad
 x_s=\frac{s-4m^2}{\Delta_c},\quad r_\lambda=\lambda_{\rm rep}\Delta_c/g_c^2,\\
 \lambda_{\rm eff}&=\frac{g_c^2}{\Delta_c}
 \left(r_\lambda-\frac{x_s}{1-x_s}\right),\quad
 \frac{\sigma(s)}{\sigma_0}=\left(\frac{\lambda_{\rm eff}}{\lambda_{\rm rep}}\right)^2
 \frac{4m^2}{s}.
\end{align}
A physical resonance needs a width and inclusive rates without double-counting
real intermediate $X$. The existing screening example is $m=1$ MeV,
$u_*=.02$, $r_\lambda=2$, $\sigma_0/m=.3\ \mathrm{cm^2/g}$,
$g_c\simeq2168$ eV, $\lambda_{\rm tot}\simeq.00881$.
It is a selected benchmark, not an identified particle.

\section{A distinct microscopic response candidate}
The separate nodal-fermion model has
\begin{align}
 \mathcal L_{\varphi P}&=\frac{Z_\varphi}{2}\dot\varphi^2
 +\frac{I_P}{2}|\dot{\mathbf P}|^2
 -\frac C2|\nabla\varphi-\mathbf P/C|^2-q_b\rho_b\varphi,\\
 H_F&=\frac{k_x^2+k_y^2-k_0^2}{2m_*}\Gamma_1+vk_z\Gamma_2
 +h\sum_{j=1}^3P_j\Gamma_{j+2},\quad
 \{\Gamma_i,\Gamma_j\}=2\delta_{ij}.
\end{align}
For $N$ occupied negative bands, zero temperature and zero chemical potential,
let $\Delta_f=h|\mathbf P|$ and $A=Nm_*/(2\pi v)$. With an energy-tube
cutoff $\Lambda_f$,
\begin{equation}
 V_F(\Delta_f)-V_F(0)=-\frac A3
 [(\Lambda_f^2+\Delta_f^2)^{3/2}-\Lambda_f^3-|\Delta_f|^3].
\end{equation}
Subtracting the analytic quadratic term gives leading
$V_{F,\rm ren}=\beta_3|\mathbf P|^3/3$, $\beta_3=Ah^3>0$.
The local cubic approximation gives
\begin{equation}
 V_\infty^4=\mathcal K\Mb,\qquad
 \mathcal K=\frac{q_b^3}{4\pi C^3\beta_3}.
\end{equation}
The required $\mathcal K=Ga_0$ remains an unfixed calibration.
In the separately derived gapless spatial regime,
\begin{equation}
 \mathbf D/C=\eta|\nabla|\mathbf P,\quad
 \eta=Nm_*h^2\pi/32,\quad
 V_\infty^2=\frac{q_b^2\Mb}{2\pi^2 C^2\eta},
\end{equation}
so $V^4\propto\Mb^2$, not the target exponent. The crossover is unsolved.
Finite temperature or occupancy further changes the local response:
\begin{equation}
 \partial_{\Delta_f}V_F=2AT\Delta_f\ln[2\cosh(\Delta_f/(2T))],\qquad
 \left.\partial_{\Delta_f}V_F\right|_{T=0}
 =A\Delta_f\max(\Delta_f,|\mu|).
\end{equation}
For $\Delta_f\ll T$, the leading potential is quadratic. This already-known
obstruction is not a new route discovered here. Positive homogeneous
response/absorption does not establish galaxy-age settling or covariant lensing.
No interaction connects this fermion medium to $\Phi_c,X$ in the current model.

\section{Route A: a Standard Model production bridge, actually screened}
\textbf{Proposal, not an adopted replacement.} Keep $\Phi_c,X$ as Standard
Model singlets with their global charges, and introduce
\begin{equation}
 \mathcal L_{\rm portal}=-\lambda_{HX}|H_{\rm SM}|^2|X|^2
 -\lambda_{H\Phi}|H_{\rm SM}|^2|\Phi_c|^2.
\end{equation}
This preserves the Standard Model gauge group and the dark $U(1)$. Initially
take $\lambda_{H\Phi}=0$ to screen the simplest extension. The broken-phase
channel is $h\to XX^*$ followed by $X\to\Phi_c\Phi_c$ and its conjugate.
It creates four daughter quanta, two of each sign, with net charge zero.
No net charge asymmetry is generated. In that neutral state the bound
$E\ge m|Q|$ supplies no nonzero abundance bound.

Freeze-in is established prior art \cite{hall}; this application supplies
the missing source interaction rather than deriving the Standard Model.
For a Maxwell--Boltzmann thermal Higgs bath, negligible inverse reactions,
radiation domination, constant $g_s,g_\rho$, high enough reheating, unit
cascade branching, and no number-changing dark reprocessing, integration
of the decay collision term gives
\begin{align}
 \dot n_\Phi+3Hn_\Phi&=4\gamma_{h\to XX^*},\\
 \gamma_{h\to XX^*}&=\frac{m_h^2T\Gamma_h}{2\pi^2}K_1(m_h/T),\\
 Y_\Phi&=4\frac{135}{8\pi^3(1.66)g_s\sqrt{g_\rho}}
 \frac{M_\Pl\Gamma_h}{m_h^2},\\
 \Gamma_h&=\frac{\lambda_{HX}^2 v_H^2}{16\pi m_h}
 \sqrt{1-4M^2/m_h^2}.
\end{align}
Here $M_\Pl=G^{-1/2}$, $Y_\Phi$ counts particles plus antiparticles, and
$Y_\Phi=2Y_{X+X^*}$. These assumptions define a production control, not
a solved dark-sector cosmology. With stipulated $\Omega_c h^2=.12$,
the illustrative conversion $\Omega_c h^2=2.742\,10^8(m/\mathrm{GeV})Y_\Phi$
and $m_h=125$ GeV, $v_H=246$ GeV, $g_s=g_\rho=106.75$, the control gives
\begin{equation}
 Y_\Phi=4.3764\,10^{-7},\quad
 \Gamma_h=4.7104\,10^{-19}\ \mathrm{GeV},\quad
 \lambda_{HX}=2.2115\,10^{-10}.
\end{equation}
The coupling is calibrated to abundance; abundance is not derived.
Assuming total Higgs width $4.1$ MeV gives illustrative branching
$1.15\,10^{-16}$, not a promised collider signal.

\subsection{First obstacle: physical masses and the scalar potential}
After electroweak symmetry breaking,
\begin{equation}
 M^2=M_0^2+\lambda_{HX}v_H^2/2,\quad
 m^2=m_0^2+\lambda_{H\Phi}v_H^2/2,\quad
 \Delta_c=\Delta_0+(\lambda_{HX}-4\lambda_{H\Phi})v_H^2/2.
\end{equation}
At the screened point the source shift is $6.6915\,10^{-6}$ GeV$^2$,
or $4182.2\Delta_c$. Moreover $M_0^2=-2.6899\,10^{-6}$ GeV$^2<0$.
The original carrier has no $|X|^4$ stabilizer. Holding $H=\Phi_c=0$ and
sending $|X|\to\infty$ then makes this minimally extended classical
potential unbounded. Thus this particular positive-portal, fixed-physical-mass
benchmark is not an admitted complete scalar model as written.
Forcing $M_0^2\ge0$ bounds
\begin{equation}
 \lambda_{HX}\le2M^2/v_H^2\simeq1.3225\,10^{-10},\qquad
 (\Omega_c h^2)_{\max}\simeq.0429
\end{equation}
within the same simplified production calculation.
The following subsection executes the positive-quartic repair. Thermal
masses, additional production channels, and changed mass/coupling points
remain distinct continuations. Setting $\lambda_{H\Phi}=\lambda_{HX}/4$
removes this tree-level contribution to the gap but is an unprotected
relation here, not an explanation. The dimensional dark-loop diagnostic
$g_c^2/(16\pi^2\Delta_c)=1.86\,10^{-5}$ is not a renormalized loop result.

\subsection{An explicit globally stable scalar extension}
Keep the desired physical masses in $V_c$ and set
$t=|H_{\rm SM}|^2$, $w_H=t-v_H^2/2$. Specify the new candidate
\begin{equation}
 V_{\rm joined}=V_c+\lambda_H w_H^2
 +\lambda_{HX}w_H|X|^2+\lambda_X|X|^4,
 \qquad \lambda_H=\frac{m_h^2}{2v_H^2}.
\end{equation}
This includes the compensating bare source mass and is gauge invariant.
Completion of a square gives the exact identity
\begin{equation}
 V_{\rm joined}=V_c+
 \lambda_H\left(w_H+\frac{\lambda_{HX}|X|^2}{2\lambda_H}\right)^2
 +\left(\lambda_X-\frac{\lambda_{HX}^2}{4\lambda_H}\right)|X|^4.
\end{equation}
Therefore $\lambda_X\ge\lambda_{HX}^2/(4\lambda_H)$ suffices for
$V_{\rm joined}\ge V_c\ge0$ for all scalar configurations, not just
near the chosen vacuum. The vacuum $X=\Phi_c=0$, $t=v_H^2/2$ is a global
minimum. For the benchmark the sufficient floor is
$\lambda_X\ge9.47\,10^{-20}$; choose $\lambda_X=10^{-18}$ as an explicit
new input. The negative bare source mass at $H=0$ no longer implies an
unbounded total potential. This is a classical boundedness result, not
a claim about the early-universe phase or radiative naturalness.

With canonical scalar kinetic terms, the coupled source and Higgs equations
include, before the additional Standard Model gauge/fermion contributions,
\begin{align}
 (\Box+M^2)X+g_c\Phi_c^2+
 \lambda_{HX}w_H X+2\lambda_X|X|^2X&=0,\\
 D_\mu D^\mu H_{\rm SM}+
 (2\lambda_Hw_H+\lambda_{HX}|X|^2)H_{\rm SM}&=0.
\end{align}
The daughter equation is unchanged at tree level in this specified
$\lambda_{H\Phi}=0$ model. Both portal forces come from the same potential;
the scalar subsystem exchanges energy reciprocally and retains the dark
$U(1)$. With vacuum energy subtracted, in isolated flat-space classical
scalar theory, the completed-square lower bound also retains $E\ge m|Q|$:
the charge-one kinetic/mass terms bound $m|Q_\Phi|$, the charge-two
$4m^2|X|^2$ terms bound $m|Q_X|$, and all added energy is nonnegative.
This extension is an admitted classical scalar production candidate, not
yet the polarization medium or a cosmologically viable theory. The portal
mass-gap sensitivity, quantum corrections, actual thermal Higgs phase,
and dark kinetic equations remain to be evaluated.

\subsection{Second obstacle: momentum and dark reprocessing}
In the rest frame of a Higgs, the source boost is $\gamma_X=m_h/(2M)$.
For isotropic source decay,
\begin{equation}
 E_\Phi=\gamma_X(M/2+\beta_X m u_*\cos\theta),\qquad
 \langle E_\Phi\rangle=m_h/4.
\end{equation}
The benchmark gives $\gamma_X=31243.8$ and daughter energies
$30.625$--$31.875$ GeV. Rest-frame coldness is not cosmological birth coldness.
Early redshifting could nevertheless cool these daughters; the boost alone
does not exclude the model.

For a deliberately collisionless monochromatic control produced at
$T_p=m_h/3$, visible entropy conservation gives
\begin{equation}
 a_p=\frac{T_{\gamma0}}{T_p}(g_{s0}/g_{sp})^{1/3},\quad
 a_{\rm nr}=a_p(p_p/m),\quad
 v(a)/c=\frac{a_{\rm nr}}{\sqrt{a^2+a_{\rm nr}^2}}.
\end{equation}
Approximating the entire interval as radiation dominated gives the exact
control integral
\begin{equation}
 \lambda_{\rm fs,RD}=\frac{c a_{\rm nr}}{H_0\sqrt{\Omega_r}}
 \left[\operatorname{asinh}(a_{\rm eq}/a_{\rm nr})
 -\operatorname{asinh}(a_p/a_{\rm nr})\right].
\end{equation}
With the declared standard-cosmology inputs in the script, this is
$4.3728\,10^{-4}$ Mpc and $p/m=1.9890\,10^{-7}$ at equality.
These are kinematic controls, not a transfer function or valid small-scale
structure constraint. The source thermal spectrum and changing degrees of
freedom are not evolved in this control.

It is also not self-consistent to retain collisionless evolution with the
adopted cross section. At visible $T=1$ keV, well below the source resonance,
using the total benchmark number density, isotropic equal-speed control
$\langle v_{\rm rel}\rangle=4v/3$, and applying the threshold cross section
to that mixture gives
\begin{equation}
 \frac{\Gamma_{22}}H=\frac{Y_\Phi s(T)\sigma_0\langle v_{\rm rel}\rangle}{H(T)}
 \simeq1374.
\end{equation}
The particle/antiparticle channel matrix must be derived; even retaining only
the same-sign half-density estimate leaves a rate much greater than $H$.
This rejects the collisionless continuation of this control, not cold
interacting matter. Momentum-resolved dark collisions, inverse conversions
and allowed number-changing reactions must replace it. The required equation is
\begin{equation}
 (\partial_t-Hp\partial_p)f_i(p,t)
 =C_i^{\rm SM}+C_i^{1\leftrightarrow2}+C_i^{2\leftrightarrow2}+cdots,
\end{equation}
with reciprocal collision integrals and energy/charge conservation.
The Higgs bridge now has an explicit tree-stable scalar extension. Its
collisional kinetics and cosmological phase history remain unresolved.

\section{Route B: an exact exchange--acceleration consistency test}
Define $Q$ in mass-density per time, positive from dark energy into cold
matter. With conserved baryons, constant $G,c,\kappa$, and the scale link
in equation~\eqref{eq:scale},
\begin{align}
 \dot\rho_c+3H\rho_c&=Q,\\
 \dot\rho_\DE+3H(1+w)\rho_\DE&=-Q,\\
 2\frac{\dd\ln a_0}{\dd\ln a}
 &=-3(1+w)-\frac{Q}{H\rho_\DE},\\
 \frac{\dd\ln(\rho_c/\rho_b)}{\dd\ln a}&=\frac{Q}{H\rho_c}.
\end{align}
Eliminating $Q$ gives the \textbf{DERIVED} conditional null relation
\begin{equation}
 \boxed{2\frac{\dd\ln a_0}{\dd\ln a}+3(1+w)
 +\frac{\rho_c}{\rho_\DE}
 \frac{\dd\ln(\rho_c/\rho_b)}{\dd\ln a}=0.}
\label{eq:null}
\end{equation}
Proof: differentiate $a_0^2\propto\rho_\DE$, insert the second continuity
equation, then substitute the fourth equation. No cosmological solution
or microscopic choice of $Q$ is required for that algebraic implication.
The code checks 27 rational cases as controls of the implementation.

For a vacuum-like component $w=-1$, exactly constant $a_0$ forces background
$Q=0$. A homogeneous net transfer from cold matter into vacuum raises
$\rho_\DE$ and $a_0$; transfer in the other direction lowers them.
The interacting-vacuum equations are known \cite{wands}; that source defines
$Q$ with the opposite sign. Our equation is their continuity system combined
with this programme's postulated scale link, not a new general cosmology.

Important scope: $Q$ here is a cosmic mean net exchange, not automatically
the local loss of halo kinetic energy. Halo cooling may deposit energy in
dark radiation or other excitations instead of vacuum. Spatially local
exchange, local gravitational work, and cosmological averaging need their own
covariant model. A constant $T_v^{\mu\nu}=-\epsilon_v g^{\mu\nu}$ has zero
covariant divergence. A varying vacuum can exchange energy but requires
$\nabla_\mu T_v^{\mu\nu}=-\nabla^\nu\epsilon_v$ and a compatible receiver.
For variable $G$ or $\kappa$, equation~\eqref{eq:null} acquires the negative
terms $-\dd\ln G/\dd\ln a-2\dd\ln\kappa/\dd\ln a$ on its left.

\textbf{Next discriminator:} derive the actual $Q^\mu$ of the proposed
formation receiver, determine its background average, and evolve abundance
and $a_0$ together. A measurable inconsistency in equation~\eqref{eq:null}
would reject the conjunction of the scale link and exchange assumptions,
not select a particle identity.

\section{Route C: can exterior free mass repair interior lensing?}
This tests one tempting escape from the inverse free-mass constraint:
place all additional positive matter outside the region where the baryonic
law must be preserved. In ordinary spherical weak-field mass lensing,
\begin{equation}
 \Sigma(R)=2\int_R^\infty\frac{\rho(r)r\dd r}{\sqrt{r^2-R^2}},\qquad
 \Delta\Sigma(R)=\frac2{R^2}\int_0^R\Sigma(S)S\dd S-\Sigma(R).
\end{equation}
For a thin shell of mass $M_s>0$ and radius $r_s>R$, put
$x=R/r_s$ and $s_x=\sqrt{1-x^2}$. Direct projection gives
\begin{align}
 \Sigma_s&=\frac{M_s}{2\pi r_s^2s_x},\quad
 M_{\rm proj}(<R)=M_s(1-s_x),\\
 \Delta\Sigma_s
 &=\frac{M_s}{\pi r_s^2}
 \left[\frac{1-s_x}{x^2}-\frac1{2s_x}\right]
 =-\frac{M_s x^2}{2\pi r_s^2s_x(1+s_x)^2}<0.
\end{align}
Newtonian interior acceleration from the shell is zero, but its excess
tangential shear is negative. Its convergence is positive; these are
different observables. By positive superposition, every finite nonnegative
spherical distribution supported at $r\ge r_0$ has
$\Delta\Sigma_{\rm ext}(R)\le0$ for $0<R<r_0$, whenever the integrals exist.
It is strictly negative for nonzero exterior mass at finite radii.
The small-$R$ limit is
\begin{equation}
 \Delta\Sigma_s=-\frac{M_sR^2}{8\pi r_s^4}+O(R^4).
\end{equation}
That leading expansion was already calculated in
\path{campaign_fresh_gravity/CFG3_principle.py}, lines 297--302.
The all-interior sign and positive-superposition application are formal
corollaries, not an independently novel force law.

The script checks five exact rational shell cases and three positive shell
mixtures. Therefore placing all repair mass beyond a protected radius cannot
increase a positive $\Delta\Sigma$ residual inside that radius while leaving
interior dynamics unchanged. This does not reject repair at larger projected
radii, nonspherical environments, interpenetrating matter, correlated
line-of-sight structure, or a different relativistic metric theory.
The next empirical step remains the measured baryon distribution and
isolation-selected lens--matter correlation required by the October 15 work
order. We have not rerun a lensing likelihood or changed an empirical verdict.

\section{Relativity, damping, and what remains to close}
The necessary joined action has the schematic form
\begin{equation}
 S=S_{\rm EH}+S_{\rm SM}+S_c+S_{\rm response}
 +S_{\rm bridge}+S_{\rm receiver}+S_{\rm counterterms}.
\end{equation}
This is an inventory, not a completed action. In a common $c=1$ convention,
\begin{equation}
 \nabla_\mu T_i^{\mu\nu}=Q_i^\nu,\qquad \sum_i Q_i^\nu=0,
 \qquad \nabla_\mu T_{\rm tot}^{\mu\nu}=0.
\end{equation}
A portal producing particles supplies neither the MOND constitutive law nor
its photon metric. In the usual weak-field convention,
\begin{equation}
 \dd s^2=-(1+2\Phi_g/c^2)c^2\dd t^2
 +a^2(1-2\Psi_g/c^2)\dd\mathbf x^2.
\end{equation}
Nonrelativistic gravity probes $\Phi_g$ and lensing probes $\Phi_g+\Psi_g$.
A separately added force on baryons cannot be assumed to bend light.
No same-action relativistic kinetic/interface proof is supplied here.

The existing Newtonian carrier starting point is
\begin{equation}
 \partial_t f+\mathbf v\cdot\nabla_x f
 -\nabla\Phi_g\cdot\nabla_v f=C[f],\quad
 \rho_c=m\int f\dd^3v,\quad
 \nabla^2\Phi_g=4\pi G(\rho_b+\rho_c).
\end{equation}
Elastic scattering redistributes energy; it cannot erase the total formation
energy. Radiation, escaping matter, a reciprocal bath, or phase mixing must
be distinguished. Conservative phase mixing can yield weak coarse settling
without strong convergence to a zero-energy resting minimum. No chosen
relaxation time is promoted to a derived clock.

\begin{longtable}{p{.20\textwidth}p{.36\textwidth}p{.35\textwidth}}
\toprule Obligation & What exists & First unresolved implication\\\midrule
\endhead
O0 source dictionary & Dated candidate equations and input hashes & One common positive-mass, charge and force mapping\\
O1 reciprocal action & Independent carrier, static and current constructions & Microscopic fields generating the adopted polarization law\\
O2 metric health & Scoped static hybrid/Einstein branches in prior work & Constrained dynamics, interfaces and lensing for the joined model\\
O3 formation & Prior finite-stock energy calculations & Actual receiver, rate and basin from independent starts\\
O4 partition & Tether/free prescription; ownership algorithm & Physical phase variable and reciprocal transition\\
O5 cosmology & Production control here; inherited cosmology in older tests & Stable potential, collision integrals, native transfer functions\\
O6 parameters & Fitted $\kappa$, stock, bounded $L_*$ & Mechanism or explicit independently measured input inventory\\
O7 joint tests & Candidate-specific successes and failures & Same candidate across all domains, both footings\\
O8 sealed data & Existing unchanged commitments & Verify public availability and execute frozen scoring when available\\
\bottomrule
\end{longtable}

The current fourth-round lensing calculation passes its relative criterion
but fails its absolute criterion: point-source $\chi^2=165.36/52$ nominal
degrees of freedom; even four fitted baryon radii give $75.73/48$ with
nominal $p=.006516<.01$. These are recorded diagnostics, not recomputed here.
They do not become successes by adding a mathematical stability theorem.
Likewise the prior static Einstein branch does not inherit the hybrid
model's static positivity theorem or prove dynamical stability.

\section{Research decisions and novelty boundary}
\textbf{Priority 1: the stabilized production candidate with full dark kinetics.}
Use the completed-square scalar extension above, determine its thermal
history, then solve the charged momentum distributions including
inverse conversion and the correct scattering channels. Success would be
an admitted potential and reproducible abundance/coldness; the source mass,
portal and initial history remain inputs unless independently fixed.
If the original 1-MeV benchmark is retained, use the explicit potential
repair and replace the failed collisionless continuation. A different mass
point requires new screening.

\textbf{Priority 2: a physical formation receiver.}
Derive reciprocal $Q^\mu$ and check energy, charge and momentum; if vacuum
is the receiver, impose equation~\eqref{eq:null}. If another component is
the receiver, propagate its perturbations and observational consequences.
This route remains open because no receiver has been specified by this report.

\textbf{Priority 3: one finite mass model for rotation and lensing.}
Use measured extended baryons, real-mass positivity, a counted stock,
the inverse constraint, and the exterior-shell sign. Predict the
isolation-selected cross-correlation rather than a fitted matter-auto-power
surrogate. The shell-only interior-repair construction is excluded in its
stated domain; broader finite population constructions remain open.

Particle production through a Higgs portal is known. Dipolar matter joining
dark energy, cold cosmological behavior and MOND-like polarization also
predates this programme \cite{blanchet}. The exchange null relation is a
formal combination of known continuity with this programme's scale postulate.
The shell result follows from standard projection, with its leading term
already in this repository. The contribution of this checkpoint is the
specific joint screening and its actionable constraints. No claim here is
certified globally unique or novel.

Literature discovery used targeted arXiv/primary-publisher web searches on
Higgs freeze-in, decay momenta, dipolar polarization, interacting vacuum,
and exterior-shell shear, checked on 16 October 2026. Primary full text was
read for \cite{hall,wands,blanchet}; abstracts of adjacent clockwork and
two-field portal work were discovery aids only. A cache lookup for
\texttt{arxiv:0911.1120v2} had no match. No PDF was retained here. This was
not a complete ADS/INSPIRE citation-graph or global novelty search.

\section{Verification and continuation checkpoint}
\texttt{checks.py} uses Python 3.9.6 standard library only. The fresh
\texttt{run\_003} records 56 passing bounded controls, including 27 exact
exchange identities, five exact shell identities, three shell mixtures,
two quadratures of the radiation-only distance and nine exact scalar
completion identities. Exact derivations are
given above; passing finite cases is not their universal proof.
The first run preceded the explicit negative-bare-mass check; run 002 added
that check, and run 003 executes the stable-quartic continuation. Earlier
runs are retained with their historical input hashes. Use run 003 for the
current script. Review is self-review, not an independent referee report.

The remaining goal is exactly O1--O7 above, plus the unchanged observational
commitments as data become available. Next execution starts with the
stabilized portal's thermal history and charged collision operators, or with a
specified reciprocal formation receiver. None is currently running in this
checkpoint. The routes are deferred/open, not exhausted and not marked solved.

\begin{thebibliography}{9}
\bibitem{hall} L. J. Hall, K. Jedamzik, J. March-Russell and S. M. West,
\emph{Freeze-In Production of FIMP Dark Matter}, arXiv:0911.1120v2
(8 December 2009), especially section 6.1, equations (6.6)--(6.10).
\url{https://arxiv.org/html/0911.1120v2}.
The decay-yield formula is specialized here with four final daughters.

\bibitem{wands} C. Kaeonikhom, H. Assadullahi, J. Schewtschenko and D. Wands,
\emph{Observational constraints on interacting vacuum energy with linear
interactions}, JCAP 01 (2023) 042, arXiv:2210.05363,
equation (2.7). \url{https://arxiv.org/pdf/2210.05363}.
The web PDF available on the check date was read; its revision was not pinned
locally. Only the displayed continuity equations, with reversed $Q$ sign,
are used, not its observational bounds.

\bibitem{blanchet} L. Blanchet and A. Le Tiec,
\emph{Model of Dark Matter and Dark Energy Based on Gravitational Polarization},
Phys. Rev. D 78, 024031 (2008), arXiv:0804.3518v2,
introduction and action framework.
\url{https://arxiv.org/html/0804.3518v2}.
This establishes overlap of the broad mechanism, not the validity of our action.
\end{thebibliography}
\end{document}
